\subsection{Splay-дерево: реализация insert, erase и find, связь с операцией splay, оценка времени работы.}

\subsection*{Insert }

Находим, куда нам нужно вставить x в дереве, запускаем splay от этого элемента, как это делали в обычном дереве, а потом вызываем splay от x. 

\textit{Асимптотика: Время insert ограниченно сверху временем работы splay. Амортизированное время работы: $O^*(\log(n))$}

\subsection*{Erase }

Поступаем точно так же, как делали в обычном бинарном дереве: разбиваем ситуацию на 3 случая, в зависимости от количества детей, но после того, как удаление произошло, запускаем splay от самой глубокой посещенной вершины.

\textit{Асимптотика: Время erase ограниченно сверху временем работы splay. Именно поэтому мы делаем splay от самой глубокой посещенной вершины. Амортизированное время работы: $O^*(\log(n))$}

\subsection*{Find }

Эта операция выполняется как для обычного бинарного дерева поиска, только после нее запускается операция splay от найденного элемента.

\textit{Асимптотика: как и в предыдущих случаях, время работы оценивается сверху временем работы splay, поэтому: амортизированное $O^*(log(n))$}

\pagebreak